\begin{proof}
\emph{(1).} Let $\phi : \shL \to \Sig \times \R^{k}$ be a fibrewise injective linear map. Its restriction to the unit vectors $\widetilde\Sig$ of $\shL$ (for any fibre metric) is a continuous map to $\R^{k} \setminus \{0\}$ with $\phi(-p) = -\phi(p)$, and $p \mapsto -p$ is the deck involution of $\widetilde\Sig \to \Sig$. Conversely a continuous $f : \widetilde\Sig \to \R^{k} \setminus \{0\}$ with $f(-p) = -f(p)$ extends to the fibrewise linear map $tp \mapsto tf(p)$, which is well defined and injective on each fibre. Dividing by the length gives maps to $S^{k-1}$. For $\Sig_{n}$ the line $\shL$ is $S^{2n+1} \times_{\langle J \rangle} \R$ with $J$ acting by $-1$; its unit vectors are the classes of $(v, \pm 1)$, so $\widetilde\Sig = S^{2n+1}/\langle J^{2} \rangle$, real projective space, with deck involution induced by $J$, and an odd map on it is a map $F$ on the sphere with $F(J(v)) = -F(v)$ (which then satisfies $F(-v) = F(v)$). On a definite side $k(\Sig) = e(\shL)$ by clause (1) of \claimref{synthesis.two-rungs-need-the-doubled-dying-line-stably-trivial}, which rests on clauses (1) and (2) of \claimref{synthesis.wall-rung-count-is-height-of-dying-line-class}; the walls exist by clause (3) there.

\emph{(2).} Write $w_{0} = a_{1}e_{1} + b_{1}e_{2}$, $w_{1} = a_{2}e_{3} + b_{2}e_{4}$, $w_{2} = a_{3}e_{5} + b_{3}e_{6}$, so $u = a_{2}e_{1} - b_{2}e_{2}$. For a vector $w = ae_{1} + be_{2}$, $w e_{1} w = (a^{2} - b^{2})e_{1} + 2ab\,e_{2}$, a vector of squared length $(a^{2} + b^{2})^{2}$ (the double-angle map of Theorem \claimref{synthesis.cohomology-of-the-quarter-turn-lens-spaces-from-the-equivariant-cell-structure}, clause (5)). The quarter turn sends $(a_{r}, b_{r})$ to $(-b_{r}, a_{r})$, hence $w_{r}e_{2r+1}w_{r}$ to its negative, $u$ to $-b_{2}e_{1} - a_{2}e_{2}$, $u e_{1} u$ to its negative and $w_{0}\cdot u = a_{1}a_{2} - b_{1}b_{2}$ to $-(a_{1}a_{2} - b_{1}b_{2})$. So $F(J(v)) = -F(v)$. With $A = w_{0}e_{1}w_{0}$ and $C = ue_{1}u$, $Q(A - C) = Q(w_{0})^{2} + Q(u)^{2} - 2A\cdot C$ and $(2w_{0}\cdot u)^{2} = 2Q(w_{0})Q(u) + 2A\cdot C$, because $A\cdot C = (a_{1}^{2} - b_{1}^{2})(a_{2}^{2} - b_{2}^{2}) - 4a_{1}b_{1}a_{2}b_{2}$ and $2(a_{1}a_{2} - b_{1}b_{2})^{2} - (a_{1}^{2} + b_{1}^{2})(a_{2}^{2} + b_{2}^{2}) = A\cdot C$. The three terms $A - C$, $2(w_{0}\cdot u)e_{3}$ and $w_{2}e_{5}w_{2}$ lie in the mutually orthogonal spaces of $\{e_{1}, e_{2}\}$, $\{e_{3}\}$, $\{e_{5}, e_{6}\}$, so $Q(F) = (Q(w_{0}) + Q(w_{1}))^{2} + Q(w_{2})^{2}$, which is positive for $v \neq 0$. The check verifies the oddness on a grid of $3^{6}$ points and the two norm identities on grids of $5^{4}$ and $5^{2}$ points, which decide polynomial identities of degree at most $2$ and $4$ in each variable. For general $n$ group the $n + 1$ planes in pairs, with one left over when $n$ is even, take the three-component map on each pair and the double-angle map on the single plane, in mutually orthogonal target spaces: the squared length is a sum of squares of positive definite forms in the separate groups of variables, zero only at the origin, and the number of components is $3\lfloor (n+1)/2 \rfloor + 2((n+1) \bmod 2) = \lceil 3(n+1)/2 \rceil$.

\emph{(3).} The cell complex is that of clause (2) of Theorem \claimref{synthesis.cohomology-of-the-quarter-turn-lens-spaces-from-the-equivariant-cell-structure}, and (G) is its stated background. A cochain of the cell complex of $S^{2n+1}$ with values in $\Z$ that is equivariant for the action of $g$ by $-1$ on the coefficients is determined by its value $c$ on $e_{k}$, taking $(-1)^{a}c$ on $g^{a}e_{k}$. From $\partial e_{k} = (g - 1)e_{k-1}$ ($k$ odd) and $\partial e_{k} = Ne_{k-1}$ ($k$ even) the coboundary is multiplication by $-1 - 1 = -2$ into odd degrees and by $1 - 1 + 1 - 1 = 0$ into even degrees. So the cohomology is $\mathbb{Z}_{2}$ in odd degrees, generated by the class of the cochain with $c = 1$, and $0$ in even degrees; the generator is not twice a cochain modulo coboundaries, which are the even cochains, so reduction modulo $2$ is injective. By the background (G), the obstruction to extending over the three-skeleton a map as in (1) into $\R^{3} \setminus \{0\}$, which is $1$-connected with $\pi_{2} = \Z$ twisted by the antipodal map, is a class in $H^{3}(\Sig_{2}; \Z_{x})$; it vanishes because the map $A - C + 2(w_{0}\cdot u)e_{3}$ of (2) is defined and nowhere zero on $S^{3}$, the three-skeleton of the stored cell structure. A nowhere-zero odd map into $\R^{2}$ is a nowhere-zero section of $\shL \oplus \shL$, the half-turn plane bundle $\eta^{2}$, whose Euler class $\pm 2y$ is not zero in $H^{2}(\Sig_{n}; \Z) = \mathbb{Z}_{4}$, $n \geq 1$.

\emph{(4).} Suppose $F : S^{5} \to \R^{4} \setminus \{0\}$ is continuous with $F(J(v)) = -F(v)$. Use $B$ as the complex structure of $\R^{6}$, so that $J$ is the unit of a circle group $\{t\}$ acting on $S^{5}$, with quotient the projective plane $P$; and identify $\R^{4}$ with a complex plane $W = W_{+} \oplus W_{-}$. Let $\lambda = S^{5} \times_{t} \mathbb{C}$ with $t$ acting on the fibre by $t^{4}$, and $V_{0} = S^{5} \times_{t} W$ with $t$ acting by $t^{2}$ on $W_{+}$ and $t^{-2}$ on $W_{-}$. The unit circle bundle of $\lambda$ is $S^{5} \times_{t} S^{1} = S^{5}/\{t : t^{4} = 1\} = \Sig_{2}$, with projection $\pi$ to $P$, and $\pi^{*}V_{0} = S^{5} \times_{\langle J \rangle} W$ with $J$ acting by $J^{2} = J^{-2} = -1$: the sum of four copies of $\shL$. The map $F$ is a nowhere-zero section $s$ of $\pi^{*}V_{0}$. The determinant line of $V_{0}$ is trivial (the circle acts trivially on $W_{+} \otimes W_{-}$); with a Hermitian metric, the orthogonal complement of $s$ is a line bundle isomorphic to the determinant, so trivial, and $\pi^{*}V_{0}$ is a trivial bundle. Let $D$ be the unit disc bundle of $\lambda$, $p : D \to P$, and $T = D \cup \mathrm{cone}(\Sig_{2})$ the Thom space. Glue $p^{*}V_{0}$ over $D$ to the product bundle over the cone by the trivialisation over $\Sig_{2}$: a complex plane bundle $V$ over $T$ whose restriction to the zero section is $V_{0}$. By (T3), $H^{2}(T) = \Z U$, $H^{4}(T) = \Z Uh$, $H^{6}(T) = \Z Uh^{2}$, and restriction to the zero section sends $U$ to $c_{1}(\lambda) = \pm 4h$. By (T2), $c_{1}(V_{0}) = 0$ and $c_{2}(V_{0}) = (2h)(-2h) = -4h^{2}$. Hence $c_{1}(V) = 0$ (restriction is injective in degree $2$) and $c_{2}(V) = nUh$ with $\pm 4n = -4$: $n = \mp 1$, odd. By (T1) and (T2), on Chern roots $c_{2} = \sum_{i<j}a_{i}a_{j}$ and $Sq^{2}(a_{i}a_{j}) = a_{i}^{2}a_{j} + a_{i}a_{j}^{2}$, so $Sq^{2}c_{2} = \sum_{i \neq j}a_{i}^{2}a_{j} = c_{1}c_{2} - 3c_{3} \equiv c_{1}c_{2} + c_{3}$; for $V$, of rank $2$ with $c_{1} = 0$, $Sq^{2}c_{2}(V) = 0$. But by (T1) and (T3), $Sq^{2}(Uh) = U^{2}h + Uh^{2} = \pm 4Uh^{2} + Uh^{2} \equiv Uh^{2}$, the generator of $H^{6}(T; \mathbb{Z}_{2})$, so $Sq^{2}(nUh) \neq 0$ for $n$ odd. This contradiction shows no such $F$ exists, so $k_{+}(\Sig_{2}) \geq 5$, and with (2), $k_{+}(\Sig_{2}) = 5$. For $n \geq 2$ a map $F$ on $S^{2n+1}$ restricts to the sphere of the first three planes, which $J$ preserves, so $k_{+}(\Sig_{n}) \geq k_{+}(\Sig_{2}) = 5$; with (2), $5 \leq k_{+}(\Sig_{3}) \leq 6$. In the language of obstruction theory: the primary obstruction to a map into $\R^{4} \setminus \{0\}$ is the Euler class $(2y)^{2} = 4y^{2} = 0$ of four copies of $\shL$, and what has been computed is the next one, with coefficients $\pi_{4}(S^{3}) = \mathbb{Z}_{2}$; no homotopy group is used in the proof. The check verifies the arithmetic: the ring $\mathbb{Z}_{2}[h, U]/(h^{3}, U^{2} - dhU)$, $Sq^{2}(Uh) = (d+1)Uh^{2}$, the identity $Sq^{2}c_{2} = c_{1}c_{2} + c_{3}$ on two and three roots, and, as a second route, that $nUh$ is the degree-four part of a Chern character with no part of degree six only if $n(d+1)$ is even (the Chern character of the Thom class being $U(e^{dh} - 1)/(dh)$; this route would replace (T1) by the integrality of the Chern character and is not used in the proof).

\emph{(5).} For the quotient by the group of order $d$, $\lambda$ has degree $d$, the line with holonomy $-1$ has complexification of degree $d/2$, $c_{2}(V_{0}) = -(d/2)^{2}h^{2}$ and $n d = -(d/2)^{2}$, so $n = -d/4$. For $d = 2$ there is no integral $n$: the Euler class of four copies of the tautological line on projective space is not zero (primary). For $d \equiv 4 \pmod 8$, $n$ is odd and $d + 1$ is odd: obstructed. For $d \equiv 0 \pmod 8$, $n$ is even: no obstruction at this stage. With $\beta x$ the element of order two of $H^{2} = \mathbb{Z}_{d}$, a class $a$ with $2a = \beta x$ is $\pm(d/4)y$ modulo the element of order two, and its reduction modulo $2$ is $u$ for $d \equiv 4 \pmod 8$ and $0$ for $d \equiv 0 \pmod 8$, well defined because the element of order two reduces to zero when $4$ divides $d$. The identification of the obstruction with $x \cup \rho(a)^{2}$ on general walls is not proved.
\end{proof}
